\documentclass[10pt]{article}
\usepackage[margin=0.85in]{geometry}
\usepackage{amsmath}
\usepackage{amssymb}
\usepackage{booktabs}
\usepackage{graphicx}
\usepackage{hyperref}
\usepackage{xcolor}

\IfFileExists{reports/generated/report_values.tex}{
  \input{reports/generated/report_values.tex}
}{
  \IfFileExists{generated/report_values.tex}{
    \input{generated/report_values.tex}
  }{
    \newcommand{\TrainingDecisions}{TBD}
    \newcommand{\EvaluationEpisodes}{TBD}
    \newcommand{\PPOSuccessRate}{TBD}
    \newcommand{\PPOSuccessCI}{TBD}
    \newcommand{\PPOCollisionRate}{TBD}
    \newcommand{\PPOCollisionCI}{TBD}
    \newcommand{\PPORouteCompletion}{TBD}
    \newcommand{\PPOAverageSpeed}{TBD}
    \newcommand{\PPOCompletionTime}{TBD}
    \newcommand{\StrongestBaseline}{TBD}
    \newcommand{\PairedSuccessDifference}{TBD}
    \newcommand{\PairedCollisionDifference}{TBD}
    \newcommand{\ModelSeed}{TBD}
    \newcommand{\ManifestHash}{TBD}
    \newcommand{\CARLAVersion}{0.9.16}
  }
}

\newcommand{\generatedfigure}[3]{
  \IfFileExists{reports/generated/#1}{
    \begin{figure}[ht]
      \centering
      \includegraphics[width=#2\linewidth]{reports/generated/#1}
      \caption{#3}
    \end{figure}
  }{
    \IfFileExists{generated/#1}{
      \begin{figure}[ht]
        \centering
        \includegraphics[width=#2\linewidth]{generated/#1}
        \caption{#3}
      \end{figure}
    }{}
  }
}

\title{Reinforcement Learning for Highway Tactical Decision-Making in CARLA}
\author{Portfolio Project Report}
\date{\today}

\begin{document}
\maketitle

\begin{abstract}
This project evaluates whether state-based PPO can improve highway tactical
decision-making over random, keep-lane, and transparent rule-based policies in
closed-loop CARLA evaluation. The policy observes a normalized state vector and
selects target-speed and target-lane actions, while every policy shares the same
conventional waypoint/PID controller and minimal car-following override.
Experiments use paired held-out traffic and weather scenarios. Numerical
results are inserted only from generated artifacts; until experiments run,
values remain explicitly marked TBD.
\end{abstract}

\section{Motivation and Project Question}
The project asks: \emph{Can a PPO agent learn safer and more efficient highway
overtaking and lane-change decisions than random, fixed keep-lane, and simple
rule-based policies under varying CARLA traffic conditions?} The scope is
highway tactical decision-making, not perception, end-to-end vehicle control,
or a claim of deployable safety.

\section{CARLA Highway Environment}
The Gymnasium environment uses CARLA \CARLAVersion{} and Town04 in synchronous
mode with a 0.05\,s fixed step. Highway starts are discovered from map topology:
candidates must be non-junction driving lanes, have a legal same-direction
adjacent lane, and retain sufficient forward continuity. Traffic Manager
controls seeded NPC traffic, including one structured slower lead vehicle.

\section{Markov Decision Process}
\subsection{State}
The observation is a 20-dimensional float vector. Eight ego/lane features
encode normalized speed, target speed, lateral offset, sine and cosine of
heading error, adjacent-lane availability, and lane-change activity. Four
features for each of the left, current, and right lanes encode capped front and
rear gaps and relative speeds.

\subsection{Actions and timing}
The five actions are maintain, accelerate, decelerate, change left, and change
right. A tactical action is selected every 1.0 simulated second and executed
over 20 low-level ticks. Unsafe or illegal lane-change requests execute as
maintain and are logged.

\subsection{Reward and termination}
For decision step \(t\), the deliberately compact reward is
\[
 r_t =
 w_p \frac{\Delta d_t}{d_s}
 + w_v\,\mathrm{clip}\!\left(\frac{v_t}{v_0},0,1.2\right)
 + c_{\mathrm{reject}} + c_{\mathrm{accepted}} + c_{\mathrm{terminal}}.
\]
Terminal penalties apply only to collision, sustained off-road state, or
stuck termination. Successful route-distance completion is terminal; the time
limit truncates an unfinished episode.

\section{Hierarchical Control Design}
The tactical policy specifies a target lane and speed. A speed-dependent
lookahead waypoint controller uses lateral and longitudinal PID terms,
steering-rate limiting, and mutually exclusive throttle and brake. A shared
car-following override caps speed near a lead vehicle and applies emergency
braking for a very small gap or TTC. It is identical across policies and is not
presented as a complete safety system.

\section{PPO and Baselines}
The state-based PPO policy uses a two-layer \(128\times128\) tanh MLP. Its
clipped surrogate term follows
\[
 L^{\mathrm{clip}}(\theta)=
 \mathbb{E}_t\!\left[
 \min\!\left(r_t(\theta)\hat A_t,
 \mathrm{clip}(r_t(\theta),1-\epsilon,1+\epsilon)\hat A_t\right)
 \right].
\]
The non-learning comparisons are a uniformly random tactical policy, a fixed
keep-lane policy, and a concise rule-based overtaking policy that checks the
left lane before the right lane.

\section{Experimental Protocol}
One frozen manifest pairs every policy on the same three traffic densities,
three weather presets, and ten held-out seeds. This yields 90 conditions per
policy and 360 episodes if complete. Primary outcomes include success,
collision, route completion, speed, successful completion time, lane changes,
and unsafe requests. Wilson intervals quantify binary rates; seeded
nonparametric bootstrap intervals quantify continuous and paired differences.
The model seed is \ModelSeed{}, the manifest hash is \ManifestHash{}, and the
recorded training budget is \TrainingDecisions{} decisions.

\section{Results}
The current evaluation row count is \EvaluationEpisodes{}. PPO success is
\PPOSuccessRate{} (95\% CI \PPOSuccessCI{}), collision is
\PPOCollisionRate{} (95\% CI \PPOCollisionCI{}), route completion is
\PPORouteCompletion{}, average speed is \PPOAverageSpeed{}, and successful
completion time is \PPOCompletionTime{}. The strongest baseline by observed
success is \StrongestBaseline{}; paired PPO differences are
\PairedSuccessDifference{} for success and \PairedCollisionDifference{} for
collision. TBD values mean experiments have not produced the corresponding
artifact.

\begin{center}
\IfFileExists{reports/generated/results_table.tex}{
  \input{reports/generated/results_table.tex}
}{
  \IfFileExists{generated/results_table.tex}{
    \input{generated/results_table.tex}
  }{\emph{Experiments pending.}}
}
\end{center}

\begin{center}
\IfFileExists{reports/generated/paired_table.tex}{
  \input{reports/generated/paired_table.tex}
}{
  \IfFileExists{generated/paired_table.tex}{
    \input{generated/paired_table.tex}
  }{}
}
\end{center}

\generatedfigure{training_episode_return.png}{0.9}{Training episode return and rolling mean.}
\generatedfigure{policy_success_collision.png}{0.9}{Held-out success and collision rates with confidence intervals.}
\generatedfigure{policy_route_completion.png}{0.8}{Mean route completion by policy.}

\section{Qualitative Driving Examples}
Six rendering categories cover a safe PPO overtake, a rejected unsafe PPO
request, a dense-traffic PPO success, a blocked keep-lane episode, a random
failure, and a genuine PPO failure. A category is omitted when no evaluation
episode meets its criterion; fallback selection requires explicit approval and
is recorded in the video manifest.

\section{Failure Analysis}
Failure analysis should distinguish tactical causes from controller or
simulation artifacts. Collision type, off-road state, stuck state, timeout,
front gap, minimum TTC, action counts, traffic density, and weather are retained
per episode. No successful episodes are cherry-picked for quantitative claims.

\section{Limitations}
The observation uses privileged CARLA state rather than perception. Neighbor
classification is approximate near road transitions, the low-level controller
is compact, Traffic Manager does not represent all driving behavior, and one
primary PPO seed limits training-variance claims. Closed-loop CARLA evaluation
does not imply real-world validity or a safety guarantee.

\section{Conclusion}
This repository provides a reproducible comparison of one state-based PPO
tactical policy with exactly three non-learning policies under a shared
conventional low-level controller. Conclusions will be stated only after the
paired evaluation artifacts are generated.

\end{document}
